\section{Leerdoelen}
Tijdens dit afstudertraject wordt gewerkt aan de vijf kerncompetenties van de Bachelor Engineering, namelijk: Analyseren, Professionaliseren, Realiseren, Ontwerpen en Valideren. Tijdens dit project ligt de focus voornamelijk op onderstaande leerdoelen.

\subsection{Analyseren}
De analysefase vormt de basis van dit onderzoek. Het doel is om binnen de eerste vier weken van het project de invloeden van omgevingsfactoren op de nauwkeurigheid te onderzoeken. Specifiek zal gekeken worden hoe reflecties en blokkades leiden tot onbetrouwbare data bij de huidige hardware. Daarnaast behoort bij deze competentie de theoretische verdieping in sensor fusion technieken als voorbereiding op de realisatie. Dit leerdoel wordt behaald aan de hand van een onderbouwd theoretisch kader, waarin de faalmechanismen gekwantificeerd zijn.

\subsection{Realiseren}
De kern van de opdracht is de realisatie van theorie naar een tastbaar product. Voor de competentie Realiseren staat de ontwikkeling van een functionerend 'Proof of Concept' centraal. Dit betreft een softwarematige implementatie van de gekozen sensor fusion techniek. Het doel is realistisch en haalbaar gehouden door de focus te leggen op de dataverwerking en eventuele algoritmes, en niet op de ontwikkeling van nieuwe hardware PCB's. Bovendien zal de data analyse zich alleen maar tot de 2D vlak limiteren. Dit leerdoel is succesvol afgerond wanneer het prototype in staat is om data van zowel de GNSS ontvanger als de IMU in te lezen, samen te voegen en autonoom foutieve waardes te detecteren. Dit zal tijdens de validatiefase gerealiseerd worden.

\subsection{Professionaliseren}
Deze competentie richt zich op het zelfstandig en methodisch uitvoeren van de opdracht binnen de context van MYLAPS. De focus ligt hierbij op actief projectmanagement, waaronder het controleren van de scope en het tijdig communiceren van struikelblokken. Dit kan behaald worden door frequent de progressie te rapporteren en regelmatig overlegmomenten met de bedrijfsbegeleider in te plannen. Bovendien zal structureel een logboek bijgehouden worden. Het doel is behaald wanneer er aantoonbaar met feedback is omgegaan en de communicatie aansluit bij de standaarden van de opdrachtgever.






\newpage